\chapter{Respiratory Rates}
\label{ch:respiratory-rates}

\section{Introduction}
\label{sec:introduction}

Respiratory rate, defined as the number of complete breathing cycles per
minute, is an important physiological vital sign
\cite{ala2025}. Accurate measurement of respiratory rate can provide useful
information about respiratory health and can be particularly important in
the assessment of respiratory illness. For example, pneumonia remains a
major cause of mortality among young children worldwide
\cite{who2022}. Accurate respiratory-rate measurement
is particularly valuable when direct respiratory monitoring is unavailable or
difficult to perform consistently. An automated and reproducible method for
estimating respiratory rate from another routinely measured physiological
signal may therefore be useful in both clinical and research settings.

This project investigates whether respiratory information can be extracted
from electrocardiogram (ECG) recordings. An ECG records the electrical
activity of the heart over time and contains a characteristic sequence of
features associated with each heartbeat. Of particular interest is the QRS
complex and its prominent R peak. Once successive R peaks have been
identified, the time interval between them can be calculated.

Breathing can influence heart rate through respiratory sinus arrhythmia.
During a breathing cycle, heart rate may increase during inspiration and
decrease during expiration. As a result, the time intervals between
successive R peaks may contain periodic variation related to respiration.
This provides a possible route for estimating respiratory rate indirectly
from ECG data.

The main research question addressed by this project is:

\begin{quote}
\textit{Can respiratory rate be estimated from variations in the R--R
intervals of long-duration ECG recordings?}
\end{quote}

The project also considers several related questions:

\begin{itemize}
    \item How reliably can R peaks be detected from ECG recordings with
    varying morphology and noise?

    \item How should the R--R intervals be represented and processed before
    respiratory analysis?

    \item Can frequency-domain or time-series methods identify a
    respiratory-related periodic component in the R--R interval sequence?

    \item How stable is the resulting respiratory-rate estimate over time?

    \item Are there observable differences between normal and
    apnea-labelled intervals?
\end{itemize}

The primary goal is to construct a transparent and reproducible computational
pipeline from the raw ECG waveform to a respiratory-rate estimate. Automated
classification of sleep apnea is treated as a possible extension rather than
the main objective of the project.


\section{Background and Theory}
\label{sec:background}

\subsection{Electrocardiogram and the QRS Complex}
\label{subsec:ecg}

An electrocardiogram measures electrical activity generated by the heart.
A typical cardiac cycle contains a P wave, QRS complex, and T wave. The QRS
complex is associated with ventricular depolarisation and usually contains a
prominent R peak. Because the R peak is often the most easily identifiable
feature of the ECG waveform, it provides a useful timing marker for
successive heartbeats.

If the detected R-peak times are written as

\begin{equation}
t_1, t_2, \ldots, t_N,
\end{equation}

then the corresponding R--R intervals are

\begin{equation}
RR_i = t_{i+1} - t_i.
\label{eq:rr_interval}
\end{equation}

The R--R interval is therefore the time between two consecutive heartbeats.
Changes in this interval reflect changes in instantaneous heart rate.

A major challenge is that ECG morphology can vary both between subjects and
within the same recording. Electrode placement, movement artefacts, noise,
ectopic beats, and physiological variation can all affect the waveform.
Consequently, a peak-detection method based on a fixed absolute threshold may
perform well for one record but fail for another.


\subsection{Respiratory Sinus Arrhythmia}
\label{subsec:rsa}

Respiratory sinus arrhythmia describes a relationship between respiration and
heart rate. In many individuals, heart rate increases slightly during
inspiration and decreases during expiration. This causes the R--R intervals
to vary over the breathing cycle.

If this respiratory modulation is sufficiently strong, the R--R interval
sequence can be treated as an ECG-derived respiratory signal. A periodic
component in this signal may then provide an estimate of respiratory rate.

However, R--R variability is not caused only by respiration. Other
physiological processes, measurement noise, missed R peaks, falsely detected
peaks, and ectopic beats can also modify the interval sequence. The central
challenge of the project is therefore not simply to find a periodic
frequency, but to determine whether a detected periodic component is
plausibly associated with respiration.


\subsection{Frequency-Domain Analysis}
\label{subsec:fft}

One possible method for estimating the dominant periodic component is the
discrete Fourier transform. For a uniformly sampled time series
$x_0, x_1, \ldots, x_{N-1}$, the discrete Fourier transform is

\begin{equation}
X_k =
\sum_{n=0}^{N-1}
x_n
\exp\left(
-\frac{2\pi i kn}{N}
\right).
\label{eq:dft}
\end{equation}

If the sampling frequency of the processed signal is $f_s$, then the
frequency corresponding to Fourier coefficient $k$ is

\begin{equation}
f_k = \frac{k f_s}{N}.
\label{eq:frequency_bins}
\end{equation}

A dominant spectral peak within a physiologically plausible respiratory
frequency range may provide an estimate of breathing frequency,
$f_{\mathrm{resp}}$. This can be converted from hertz to breaths per minute
using

\begin{equation}
R_{\mathrm{resp}}
=
60 f_{\mathrm{resp}},
\label{eq:resp_rate}
\end{equation}

where $R_{\mathrm{resp}}$ is the estimated respiratory rate in breaths per
minute.

A complication is that R--R intervals are naturally irregularly sampled
because their timestamps are determined by the heartbeats themselves.
Therefore, direct application of a standard FFT may require interpolation
onto a uniformly sampled time grid. The effects of this preprocessing must
be considered when interpreting the resulting spectrum.


\section{Data}
\label{sec:data}

The project uses data from the publicly available Apnea-ECG database
\cite{penzel2000}. The dataset contains long-duration single-channel ECG
recordings collected as part of a sleep-apnea study.

Several file types are relevant to the analysis:

\begin{table}[h!]
\centering
\begin{tabular}{ll}
\hline
\textbf{File type} & \textbf{Purpose} \\
\hline
\texttt{.dat} & ECG waveform data \\
\texttt{.hea} & Record metadata and signal information \\
\texttt{.qrs} & QRS-complex annotations \\
\texttt{.apn} & Sleep-apnea annotations \\
\hline
\end{tabular}
\caption{Main file types used from the Apnea-ECG dataset.}
\label{tab:file_types}
\end{table}

The waveform and header files provide the ECG signal and the information
required to interpret it. Where available, the QRS annotations provide
reference locations for QRS complexes, while the apnea annotations identify
periods labelled as normal or apneic.

The WFDB annotation format is more complex than a simple text format and
contains compact binary annotation records. The WFDB specification
\cite{physionet_wfdb} is therefore used as a reference when reading or
interpreting annotation files.

Not every record contains the same set of annotations, and the available
QRS annotations should not automatically be treated as perfectly accurate.
They are instead used as a useful reference for testing the ECG-processing
pipeline.


\section{Methods}
\label{sec:methods}

The planned analysis follows the computational pipeline

\begin{equation}
\text{ECG}
\rightarrow
\text{R peaks}
\rightarrow
\text{R--R intervals}
\rightarrow
\text{respiratory-related signal}
\rightarrow
\text{respiratory-rate estimate}.
\label{eq:pipeline}
\end{equation}


\subsection{ECG Loading and Visualisation}
\label{subsec:loading}

The first stage is to load the ECG waveform and associated metadata and
convert the signal from sample number to physical time. Representative
segments of the ECG are then visualised together with the available QRS and
apnea annotations.

This stage serves two purposes. First, it verifies that the data and
annotations have been interpreted correctly. Second, visual inspection helps
identify differences in morphology, noise level, and signal amplitude that
may affect later R-peak detection.

The initial objective is to produce a figure containing a short ECG interval,
the corresponding QRS locations, and any overlapping apnea annotation. This
provides a baseline visualisation against which the later automated analysis
can be checked.


\subsection{R-Peak Detection}
\label{subsec:r_peak}

The next stage is to identify the R peaks directly from the ECG waveform.
The detector must be sufficiently robust to deal with changes in signal
amplitude and morphology.

A simple fixed-amplitude threshold is unlikely to generalise across all
records. The detection method will therefore use relative signal properties
such as local peak prominence, minimum peak separation, or an adaptive
threshold. A minimum temporal separation between candidate peaks can also be
used to prevent multiple detections within a single QRS complex.

Where QRS annotations are available, detected R peaks will be compared with
the annotation locations. The comparison will be used to identify common
failure modes such as missed peaks, false peaks, and timing offsets.

Previous work has shown that the choice of QRS-detection algorithm can affect
derived heart-rate variability measurements, making this validation an
important part of the analysis \cite{weinstein2026}.


\subsection{R--R Interval Construction}
\label{subsec:rr}

After detecting the R peaks, the intervals are calculated using
Equation~\ref{eq:rr_interval}. Each interval is associated with a
representative time, for example the timestamp of the later R peak or the
midpoint between the two peaks.

Before respiratory analysis, implausible R--R intervals must be identified.
Abnormally short or long intervals may be caused by missed detections, false
detections, ectopic beats, or noise. Such intervals can strongly influence a
frequency-domain analysis.

The resulting R--R sequence is irregularly sampled. If a conventional FFT is
used, the cleaned interval series will therefore be interpolated onto a
regular time grid. The interpolation frequency must be high enough to capture
respiratory variation without unnecessarily increasing the amount of data.


\subsection{Respiratory-Rate Estimation}
\label{subsec:resp_estimation}

The respiratory-rate estimation stage will analyse periodic variation in the
processed R--R interval series.

The baseline approach is:

\begin{enumerate}
    \item detect R peaks from the ECG;
    \item calculate the R--R intervals;
    \item remove or correct clearly implausible intervals;
    \item interpolate the R--R sequence onto a regular time grid;
    \item remove slow trends where necessary;
    \item divide the signal into analysis windows;
    \item calculate the frequency spectrum within each window;
    \item identify the dominant plausible respiratory frequency;
    \item convert the frequency to breaths per minute using
    Equation~\ref{eq:resp_rate}.
\end{enumerate}

Using multiple time windows rather than a single spectrum of the entire
recording should make it possible to investigate how the estimated
respiratory rate changes over time.

The choice of analysis-window length introduces a time-frequency trade-off.
Longer windows improve frequency resolution but reduce the ability to track
rapid changes. Shorter windows provide better temporal localisation but
produce less precise frequency estimates.


\section{Preliminary Results}
\label{sec:results}

At the current draft stage, the project is focused on establishing and
validating the data-processing pipeline before drawing conclusions about
respiratory rate.

The first milestone is successful loading and visualisation of representative
ECG records together with their available QRS and apnea annotations. This
provides a direct check that the waveform timing and annotation timing are
consistent.

The next preliminary result will be a comparison between automatically
detected R peaks and the provided QRS annotations. This comparison will be
used to determine whether the peak-detection method is sufficiently reliable
for construction of the R--R interval series.

Once the R--R series has been validated, representative sections of the
interval sequence will be examined in both the time and frequency domains.
The main quantity of interest will be the presence and stability of a
dominant respiratory-related periodic component.

% Example figure block to activate once the corresponding figure is ready:
%
% \begin{figure}[h!]
% \centering
% \includegraphics[width=0.85\linewidth]{figures/ecg_example.pdf}
% \caption{Representative ECG segment with QRS annotations and an
% apnea-labelled interval.}
% \label{fig:ecg_example}
% \end{figure}


\section{Discussion and Current Limitations}
\label{sec:discussion}

Several limitations must be considered when interpreting ECG-derived
respiratory estimates.

First, respiratory sinus arrhythmia is not the only source of heart-rate
variability. Changes in autonomic activity, ectopic beats, movement
artefacts, and other physiological processes may also affect the R--R
intervals.

Second, errors in R-peak detection propagate directly into the R--R interval
series. A single missed peak may generate an interval approximately twice as
long as expected, while an additional false peak may create two artificially
short intervals. Reliable peak detection and outlier handling are therefore
essential.

Third, interpolation changes the original irregular R--R data and may
introduce additional assumptions into the analysis. The interpolation method
and resampling rate should therefore be chosen carefully.

A further limitation is the absence of a continuous direct respiratory-rate
reference signal in the current project setup. The apnea annotations indicate
periods of sleep apnea, but they are not equivalent to point-by-point
ground-truth respiratory-rate measurements. Consequently, a spectral peak
that appears physiologically plausible should not automatically be interpreted
as a fully validated respiratory-rate measurement.

Finally, ECG morphology varies between subjects and recording conditions.
A method that performs well on a single record may therefore not generalise
to the full dataset.


\section{Planned Next Steps}
\label{sec:next_steps}

The next stages of the project are:

\begin{enumerate}
    \item complete and validate the ECG and annotation loading pipeline;

    \item implement robust R-peak detection and compare the results with
    available QRS annotations;

    \item construct and clean the R--R interval time series;

    \item investigate interpolation and preprocessing choices for the
    irregularly sampled R--R sequence;

    \item implement a baseline spectral respiratory-rate estimator;

    \item evaluate the stability of respiratory-rate estimates across
    multiple time windows and ECG records;

    \item compare respiratory-related features between normal and
    apnea-labelled intervals;

    \item if time permits, investigate whether simple statistical features
    can help distinguish normal and apnea-labelled intervals.
\end{enumerate}

The priority remains the development of a reliable and interpretable
respiratory-rate estimation pipeline. More complex classification or
machine-learning methods will only be considered after the baseline analysis
has been validated.


\section{Conclusion}
\label{sec:conclusion}

This project investigates whether respiratory information can be extracted
indirectly from long-duration ECG recordings. The proposed approach uses the
timing of R peaks to construct an R--R interval time series and then searches
for periodic variation associated with respiration.

The central computational challenge is to separate respiratory-related
variation from errors in peak detection, normal heart-rate variability,
artefacts, and other physiological effects. The draft therefore focuses on
building and validating each stage of the processing pipeline before
interpreting the final respiratory-rate estimate.

The next phase of the project will determine whether the resulting
R--R-derived signal contains a sufficiently stable respiratory component to
support meaningful estimates over time and whether these features differ
between normal and apnea-labelled intervals.
